\begin{tabular}{lrrrr}
\toprule
\textbf{Punto} & \textbf{Desacuerdo [\%]} & \textbf{SD [keV]} & \textbf{Similitud coseno} & \textbf{$R^2$} \\
\midrule
A @ 1600 V & $+9.7$ & 36.1 & 1.00000000 & 0.9713 \\
A @ 1650 V & $+12.6$ & 31.0 & 1.00000000 & 0.9787 \\
A @ 1700 V & $+14.7$ & 43.3 & 1.00000000 & 0.9539 \\
B @ 1600 V & $+9.0$ & 22.5 & 1.00000000 & 0.9896 \\
B @ 1650 V & $+8.8$ & 20.3 & 1.00000000 & 0.9917 \\
B @ 1700 V & $+9.9$ & 23.9 & 1.00000000 & 0.9882 \\
C @ 1600 V & $-11.3$ & 49.0 & 1.00000000 & 0.9716 \\
C @ 1650 V & $-8.0$ & 36.4 & 1.00000000 & 0.9829 \\
C @ 1700 V & $-9.1$ & 53.1 & 1.00000000 & 0.9663 \\
D @ 1600 V & $+19.3$ & 49.3 & 1.00000000 & 0.9360 \\
D @ 1650 V & $+16.5$ & 47.2 & 1.00000000 & 0.9427 \\
D @ 1700 V & $+2.2$ & 0.6 & 1.00000000 & 1.0000 \\
\bottomrule
\end{tabular}
\par\smallskip
\begin{minipage}{0.95\linewidth}\footnotesize Las dos ultimas columnas son las metricas con las que el paper valida su calibracion (su Fig. 3). La similitud coseno da 1 EXACTAMENTE en todos los puntos, incluido el que discrepa mas: entre dos rectas por el origen es 1 sean cuales sean las pendientes, asi que no puede ver este modo de falla. El $R^2$ se calcula sobre dos puntos y con la recta forzada por el origen; con ordenada libre y dos puntos daria 1 exacto.\end{minipage}
